设弱激发荧光显微镜中的有效两能级分子自发辐射率为 \(\Gamma>0\)。激发率为 \(W_{\rm ex}(r)=W_0e^{-r^2/(2\sigma^2)}\)，其中 \(0<W_0\ll\Gamma\)；受激耗尽光在研究区域使用二次零心模型 \(W_{\rm dep}(r)=\Gamma\beta r^2/w^2\)，其中 \(\beta,w>0\)。此二次模型被指定为全平面数学延拓，用来比较径向亮度和二维总光子数，忽略光漂白、真实光斑有限孔径及复杂能级。令 \(\alpha=1/(2\sigma^2)\)、\(c=\beta/w^2\)。
(1) 建立稳态激发态布居的速率方程并推导弱激发荧光 \(F(r)\) 及中心归一化点扩散函数；解释受激耗尽的物理作用和弱激发近似的必要性。
(2) 从中心邻域的对数强度展开求有效高斯宽度 \(\sigma_{\rm eff}\)，再不依赖小 \(r\) 展开，利用Lambert \(W_0\) 函数精确解半高位置 \(r_{1/2}\)，证明唯一性和无耗尽极限。
(3) 对所给全平面模型精确积分 \(N(c)=2\pi\int_0^\infty rF(r)dr\)；用指数积分 \(E_1\) 表示 \(N(c)/N(0)\)，证明增强耗尽使总光子数单调减少并求大 \(c\) 对数渐近。
(4) 在给定曝光时间和入射激发率下，比较大 \(c\) 时的局部宽度缩减与Poisson散粒噪声代价，指出为何仅有PSF缩窄不能直接保证实验定位误差无限下降；需要明确辨析有限光子预算和未建模噪声源。